\documentclass{beamer}

\usepackage[utf8]{inputenc}
\input{../helper}

\title{Introdução à Computação \\ lendo e escrevendo em arquivos}
\author{Alexandre Rademaker}
\date{}

\begin{document}

\begin{frame}
 \maketitle
\end{frame}


\begin{frame}[allowframebreaks,fragile]{scanf}

  \ilcode{scanf} recebe uma string template que contém designadores de
  tipos como \ilcode{\%i}, então a entrada padrão (STDIN) é lida até
  que o padrão seja reconhecido.

  Par cada designador mencionado no template, um parâmetro adicional
  deve ser passado, um endereço para um tipo compatível com o
  designador.

  \ilcode{scanf} deverá efetivamente modificar o endereço indicado,
  preenchendo o valor lido, ou seja precisamos passar uma
  'referência'.

  \framebreak

  \lstinputlisting[style=tinyc,caption=scanf1.c]{src/scanf1.c}

  Veja o \ilcode{bus error}! Não alocamos memória para armazenar o
  valor lido. Mas como fazer isso? Qual tamanho devemos alocar?

  \framebreak

  \lstinputlisting[style=tinyc,caption=src/scanf2.c]{src/scanf2.c}

  \ldots ou seja, usar \ilcode{scanf} para ler strings pode ser
  perigoso! 

  \framebreak

  \lstinputlisting[style=tinyc,caption=src/scanf3.c]{src/scanf3.c}

  e ainda existem \ilcode{fgets} e \ilcode{gets}.  No MacOS ou Linux,
  use as \textbf{man pages} com \ilcode{man gets}.

\end{frame}  


\begin{frame}[allowframebreaks,fragile]{lendo e escrevendo em arquivos}

  Todas as principais funções para ler e escrever em arquivos ficam na
  biblioteca \ilcode{stdio.h}.

  As funções mais comuns são: \ilcode{fopen}, \ilcode{fclose},
  \ilcode{fgetc}, \ilcode{fputc}, \ilcode{fread}, \ilcode{fwrite},
  \ilcode{fscanf}, \ilcode{fprintf}.

  \framebreak

  Veja arquivo phonebook.c
  % \lstinputlisting[style=tinyc,caption=src/phonebook.c]{src/phonebook.c}

  A função \ilcode{fopen}, o tipo \ilcode{FILE} e \ilcode{fprintf} é a
  versão de \ilcode{printf} que escreve em um arquivo.

  \framebreak

  antes

  \begin{lstlisting}[style=code,caption=phonebook.csv,autogobble]
  name,number
  \end{lstlisting}

  depois

  \begin{lstlisting}[style=code,caption=phonebook.csv,autogobble]
    name,number
    Alexandre Rademaker,+55 21 12345
  \end{lstlisting}

  \framebreak

  Podemos ler arquivos \href{https://www.file-recovery.com/jpg-signature-format.htm}{jpeg}

  \lstinputlisting[style=tinyc,caption=jpeg.c,autogobble]{src/jpeg.c}

  \framebreak

  O formato \href{https://en.wikipedia.org/wiki/BMP_file_format}{BMP}
  é outro formato popular para armazenar imagens. 

  \begin{lstlisting}[style=tinyc,autogobble,caption=src/filter]
    #include "helpers.h"
  
    void filter(int height, int width, RGBTRIPLE image[height][width])
    {
      // Loop over all pixels
      for (int i = 0; i < height; i++)
      {
        for (int j = 0; j < width; j++)
        {
          image[i][j].rgbtBlue = 0x00;
          image[i][j].rgbtGreen = 0x00;
        }
      }
    }
  \end{lstlisting}

  A função acima percorre todos os pixels em um array bidimensional
  e define os valores de azul e verde como 0.
  
\end{frame}

\end{document}

%%% Local Variables:
%%% mode: latex
%%% TeX-master: t
%%% End:
